% =============================================================================
\section{Experimental Setup}
\label{sec:experiments}
% =============================================================================

\subsection{Datasets and Metrics}

We evaluate on two groups of benchmarks:

\textbf{Primary benchmarks} (TREC Deep Learning):
\begin{itemize}[leftmargin=*]
\item \textbf{TREC DL 2019}~\citep{craswell2020overview}: 43 queries with graded relevance judgments (0--3) over the MS MARCO passage corpus (8.8M passages).
\item \textbf{TREC DL 2020}~\citep{craswell2021overview}: 54 queries with the same passage corpus and grading scheme.
\end{itemize}

\textbf{Generalizability benchmarks} (BEIR, in progress): To assess domain generalizability, we additionally run a representative subset of BEIR~\citep{thakur2021beir} datasets on three models (Flan-T5-XL, Qwen3-8B, Qwen3.5-9B). These runs are still completing on a remote cluster and will be folded into the camera-ready version. The TREC DL 2019/2020 results below are the basis of all of our headline claims.

For all datasets, we rerank the top-100 BM25 results from Pyserini~\citep{lin2021pyserini} using default parameters.

\textbf{Effectiveness metrics}: NDCG@10 is our primary metric, following prior work on LLM-based reranking~\citep{zhuang2024setwise, sun2023chatgpt, qin2024large}. We also report MAP@100 for completeness in the main results table.

\textbf{Efficiency metrics}: We report (i) the number of LLM comparison calls per query (\#Comps), (ii) total prompt tokens per query, (iii) total completion tokens per query, and (iv) wall-clock time per query.

\textbf{Statistical testing}: We compare each challenger family (DualEnd, BottomUp, BiDir) against the best TopDown baseline for the same model-dataset configuration using two-sided paired approximate randomization with $100{,}000$ samples. Mean-delta uncertainty is reported as a paired bootstrap 95\% confidence interval ($20{,}000$ resamples). Multiple testing is controlled by Bonferroni correction within each challenger family across the 18 TREC DL configurations. Per-query \texttt{ndcg\_cut\_10} values are read directly from the saved \texttt{.eval} files in the experiment output tree, so the test statistic is grounded in the same numbers as the main results table. Full per-config tables and verdicts appear in the supplementary \texttt{SIGNIFICANCE\_TESTS} artifact and are summarized in \S\ref{sec:rq2}.

\subsection{Models}

We select models from three architecture families spanning a range of sizes to assess generalizability:

\textbf{Encoder-decoder (Flan-T5)}: Instruction-tuned sequence-to-sequence models that natively support both generation and likelihood scoring modes via direct logits over single-character labels.
\begin{itemize}[leftmargin=*]
\item \textbf{Flan-T5-Large} (780M parameters)
\item \textbf{Flan-T5-XL} (3B parameters) --- used as the primary model for the bulk of ablations.
\item \textbf{Flan-T5-XXL} (11B parameters) --- for scaling analysis.
\end{itemize}

\textbf{Decoder-only (Qwen3)}: Instruction-tuned causal language models with chain-of-thought reasoning. These models emit \texttt{<think>...</think>} reasoning blocks, which we strip during output parsing. The main paper tables use generation mode; the codebase also supports a likelihood follow-up that scores short teacher-forced answer strings such as \texttt{Passage A} without decoding.
\begin{itemize}[leftmargin=*]
\item \textbf{Qwen3-4B} (4B parameters)
\item \textbf{Qwen3-8B} (8B parameters)
\item \textbf{Qwen3-14B} (14B parameters)
\end{itemize}

\textbf{Hybrid causal (Qwen3.5)}: A novel architecture combining Gated DeltaNet linear attention (75\% of layers) with standard softmax attention (25\%), resulting in a decoder-only model with sub-quadratic complexity. Despite the \texttt{ForConditionalGeneration} class name (which refers to its multimodal capability), these are causal language models loaded via \texttt{AutoModelForCausalLM}. The main paper tables use generation mode; the codebase also supports the same teacher-forced likelihood follow-up used for Qwen3.
\begin{itemize}[leftmargin=*]
\item \textbf{Qwen3.5-4B} (4B parameters)
\item \textbf{Qwen3.5-9B} (9B parameters)
\item \textbf{Qwen3.5-27B} (27B parameters)
\end{itemize}

All models are run in half precision (FP16/BF16) on a single NVIDIA H100 80GB GPU. Flan-T5 models use passage truncation length of 128 tokens (following \citet{zhuang2024setwise}, constrained by 512-token model context). Qwen3/3.5 models use 512 tokens (32k+ context window). We study the effect of passage length in an ablation (\S\ref{sec:ablation}).

\subsection{Configurations}

Table~\ref{tab:configs} summarizes all experimental configurations. For all setwise methods, we set \texttt{num\_child}$=3$ (4-way comparison) following~\citet{zhuang2024setwise}. We rerank the top 100 BM25 results and extract the top $k=10$. Passage truncation length varies by model family: 128 tokens for Flan-T5 and 512 tokens for Qwen3/3.5.

\begin{table}[!htbp]
\caption{Experimental configurations. All use $k{=}10$, hits${}={100}$, $c{=}3$.}
\label{tab:configs}
\small
\resizebox{\columnwidth}{!}{%
\begin{tabular}{llll}
\toprule
\textbf{Method} & \textbf{Direction} & \textbf{Algorithm} & \textbf{Extra} \\
\midrule
\topdown{}-Heap & topdown & heapsort & --- \\
\topdown{}-Bubble & topdown & bubblesort & --- \\
\bottomup{}-Heap & bottomup & heapsort & --- \\
\bottomup{}-Bubble & bottomup & bubblesort & --- \\
\dualend{}-Cocktail & dualend & bubblesort & cocktail shaker \\
\dualend{}-Selection & dualend & selection & DE selection \\
\bidir{}-RRF & bidirectional & heapsort & RRF ($k_{\text{rrf}}{=}60$) \\
\bidir{}-Weighted & bidirectional & heapsort & $\alpha{=}0.7$ \\
\bottomrule
\end{tabular}
}
\end{table}

\textbf{Implementation}: All methods are implemented in the \texttt{llm-rankers} framework~\citep{zhuang2024setwise}, extending the existing \texttt{SetwiseLlmRanker} class. The \texttt{BottomUpSetwiseLlmRanker}, \texttt{DualEndSetwiseLlmRanker}, and \texttt{BidirectionalEnsembleRanker} classes are added as new modules, alongside the three refinement variants \texttt{SelectiveDualEndSetwiseLlmRanker}, \texttt{BiasAwareDualEndSetwiseLlmRanker}, and \texttt{SameCallRegularizedSetwiseLlmRanker} described in \S\ref{sec:refinement}. The bidirectional ensemble uses Python's \texttt{\_\_new\_\_} constructor pattern to share model weights between the top-down and bottom-up rankers, avoiding redundant GPU memory usage. For Qwen3-family thinking models, output parsing handles \texttt{<think>...</think>} reasoning blocks via regex stripping, with \texttt{max\_new\_tokens}${}=256$ to allow the thinking block to complete before the answer (and 512 for the longer dual-output prompt). The causal likelihood follow-up scores short continuations such as \texttt{Passage A} with teacher forcing rather than assuming a single stable label token.
